% Figure 6.x: a shear has det = 1 and preserves area.
\begin{tikzpicture}[scale=1.0,>=stealth]
  % left: unit square
  \begin{scope}
    \draw[thick,fill=accentsoft] (0,0) -- (1,0) -- (1,1) -- (0,1) -- cycle;
    \draw[->,thick] (0,0) -- (1,0) node[midway,below,font=\footnotesize] {$\mathbf e_1$};
    \draw[->,thick] (0,0) -- (0,1) node[midway,left,font=\footnotesize] {$\mathbf e_2$};
    \node[font=\small] at (0.5,-0.55) {unit square, area $=1$};
  \end{scope}
  % arrow
  \draw[->,thick,accent] (1.9,0.6) -- (2.9,0.6)
    node[midway,above,font=\footnotesize] {$A=\begin{bmatrix}1&1\\0&1\end{bmatrix}$};
  % right: parallelogram
  \begin{scope}[xshift=4.6cm]
    \draw[thick,fill=accentsoft] (0,0) -- (1,0) -- (2,1) -- (1,1) -- cycle;
    \draw[->,thick] (0,0) -- (1,0) node[midway,below,font=\footnotesize] {$A\mathbf e_1$};
    \draw[->,thick] (0,0) -- (1,1) node[midway,left,font=\footnotesize] {$A\mathbf e_2$};
    \draw[dashed,gray] (1,0) -- (2,1);
    \node[font=\small] at (1.0,-0.55) {parallelogram, area $=1$};
  \end{scope}
  \node[font=\footnotesize,align=center] at (3.4,1.7) {$\det A=1$: shear slides,\\never scales area};
\end{tikzpicture}
